%% -----------------------------------------------------------
\section{Conditions for Safe, Effective Deployment}
%% -----------------------------------------------------------

Six design principles emerge from the evidence as necessary conditions for equity-centered deployment.

\subsection{Community Co-Design as Safety Imperative}

Co-design with target communities functions as both an efficacy requirement and a safety requirement. Only half of \gls{lgbtq} interventions included meaningful community input, yet those employing co-design demonstrated higher effectiveness, better engagement, and greater trust~\cite{fowler_2023}. Community involvement helped identify invisible harms, incorporate affirming language, and build the credibility that drives sustained engagement~\cite{fowler_2023}. Outsider-developed interventions provided advice immediately recognizable as culturally inappropriate by anyone with lived experience~\cite{soubutts_2024}. For neurodivergent populations, Papadopoulos~\cite{papadopoulos_2025} specifically emphasizes that co-design with neurodivergent developers and users produces qualitatively different tools from those designed by neurotypical teams alone.

\subsection{Cultural Humility Through Continuous Adaptation}

Static translation fails; culturally humble approaches require continuous community engagement. Naderbagi et al.~\cite{naderbagi_2024} established a framework distinguishing surface-level adaptation from deep-structure adaptation that addresses how populations conceptualize mental health, perceive illness causes, define treatment success, and approach help-seeking. The framework requires iterative testing with diverse users, flexibility accommodating within-group diversity, and recognition that culture extends beyond ethnicity to encompass family structure, community norms, and historical experiences of healthcare systems.

\subsection{Hybrid Human-AI Models}

Hybrid architectures balance scalability with safety. Conversational AI effectively facilitates assessments and triage when human oversight is embedded in the system architecture~\cite{rollwage_2024}. Professionally-facilitated digital programs achieve both higher engagement and better safety outcomes than fully automated systems~\cite{liu_2023}. Human escalation pathways must be built in from inception~\cite{smith_2023}. For high-risk populations, including those with active suicidal ideation or severe trauma histories, full automation represents an inadequate standard of care regardless of chatbot sophistication~\cite{de_freitas_2024,palmer_2025}.

\subsection{Transparency About Limitations}

Transparency improves outcomes and is a precondition for informed consent. Explicitly acknowledging limitations produces better results than claiming comprehensive support~\cite{drydakis_2022}. Smith et al.~\cite{smith_2023} recommend mandatory disclosure covering what conditions chatbots address, what evidence supports effectiveness, what safety protocols exist, what data privacy protections apply, and what chatbots cannot do. The American Psychological Association has reinforced this position, recommending against generic chatbot use for mental health support~\cite{apa_2025}. Palmer et al.~\cite{palmer_2025} propose a voluntary certification program that would standardize disclosure requirements across digital mental health platforms.

\subsection{Accessibility as Foundational Architecture}

Accessibility must be foundational from the earliest design stages. Khan and Seo~\cite{khan_seo_2025} establish that genuine accessibility requires involving disabled users throughout development, treating \gls{wcag} compliance as a minimum rather than a ceiling, conducting testing with actual assistive technology users, and designing interaction patterns that work across diverse ability levels by default. For neurodivergent users, accessible design additionally requires attention to cognitive load, sensory sensitivity, communication style variability, and the absence of neurotypical social cues in interaction design~\cite{papadopoulos_2025,killian_2023}.

\subsection{Fair-Aware AI Development}

Timmons et al.~\cite{timmons_2023} propose a ``fair-aware AI'' framework requiring bias assessment at every stage of the model development lifecycle, from training data curation and feature selection through model evaluation, deployment monitoring, and post-market surveillance. Their framework specifically identifies that mental health AI must be evaluated against equity outcomes (whether it reduces or widens disparities) alongside efficacy outcomes (whether it reduces symptoms)~\cite{timmons_2023}. This framework extends to racial and ethnic bias~\cite{haider_2024,omar_2025}, gender and sexuality bias~\cite{straw_2020}, and neurotypical bias~\cite{killian_2023}, providing a unified methodological approach to equity across population groups.
